\Needspace{6\baselineskip}
\begin{center}\hypertarget{native-261-article188}{}\pdfbookmark[2]{188.}{article-188}188.\end{center}
Das vorige Verfahren eignet sich vor allem für Fälle, in denen höchste
Genauigkeit erwünscht ist. Sehr oft kann man jedoch ohne Bedenken etwas
Genauigkeit preisgeben, wenn sich dadurch die Rechnung erheblich abkürzen
lässt. Das gilt besonders dann, wenn die Beobachtungen noch keinen langen
Zeitraum umfassen und daher an eine, wenn man so sagen darf, endgültige
Bahnbestimmung noch nicht zu denken ist. In solchen Fällen lässt sich
das folgende Verfahren mit beträchtlichem Gewinn anwenden.

Man wähle aus allen Beobachtungen zwei vollständige Örter $L$ und $L'$ aus
und berechne für die entsprechenden Zeitpunkte aus den Näherungselementen
die Entfernungen des Himmelskörpers von der Erde. Dann stelle man für
diese Entfernungen drei Hypothesen auf. In der ersten behalte man die
berechneten Werte bei, in der zweiten ändere man die erste Entfernung, in
der dritten die zweite. Beide Änderungen kann man nach Belieben so wählen,
dass sie der Unsicherheit entsprechen, die man bei den jeweiligen
Entfernungen noch vermutet. Unter den drei Hypothesen in der folgenden
Übersicht
\begin{center}
\begingroup
\setlength{\tabcolsep}{4pt}
\renewcommand{\arraystretch}{1.3}
\begin{tabular}{|l|c|c|c|}
\hline
 & Hyp. I & Hyp. II & Hyp. III\\
\hline
Entfernung\footnotemark\ für den ersten Ort & $D$ & $D+\delta$ & $D$\\
Entfernung für den zweiten Ort & $D'$ & $D'$ & $D'+\delta'$\\
\hline
\end{tabular}
\endgroup
\end{center}
\footnotetext{Noch bequemer ist es, anstelle der Entfernungen selbst die Logarithmen der verkürzten Entfernungen zu verwenden.}
\noindent\SourceLeaf{262}berechne man aus den beiden
Örtern $L,L'$ nach den im ersten Buch erläuterten Verfahren drei Systeme
von Elementen. Aus jedem dieser Systeme berechne man dann die geozentrischen
Örter des Himmelskörpers zu den Zeitpunkten aller übrigen Beobachtungen.
Sie seien wie folgt bezeichnet, wobei die einzelnen Längen und Breiten
beziehungsweise Rektaszensionen und Deklinationen gesondert bezeichnet werden.
\begin{center}
\begin{tabular}{@{}l@{\quad}l@{\quad}l@{\quad}l@{\quad}l@{}}
\makebox[12.8em][l]{im ersten System\hfill .\,.\,.\,.\,.\,.} & $M$, & $M'$, & $M''$ & usw.\\[.3em]
\makebox[12.8em][l]{im zweiten System\hfill .\,.\,.\,.\,.} & $M+a$, & $M'+a'$, & $M''+a''$ & usw.\\[.3em]
\makebox[12.8em][l]{im dritten System\hfill .\,.\,.\,.\,.\,.\,.} & $M+\beta$, & $M'+\beta'$, & $M''+\beta''$ & usw.
\end{tabular}
\end{center}
\noindent Ferner seien die entsprechenden beobachteten Örter
\begin{center}
\begin{tabular}{@{}l@{\quad}l@{\quad}l@{\quad}l@{\quad}l@{}}
\makebox[12.8em][l]{beobachtete Örter\hfill .\,.\,.\,.\,.\,.\,.\,.\,.\,.\,.} & $N$, & $N'$, & $N''$ & usw.
\end{tabular}
\end{center}

Soweit nun kleinen Änderungen der Entfernungen $D,D'$ proportionale
Änderungen der einzelnen Elemente und der daraus berechneten
geozentrischen Örter entsprechen, darf man Folgendes annehmen. Legt man
einem vierten System von Elementen die Entfernungen $D+x\delta,D'+y\delta'$
von der Erde zugrunde, so betragen die daraus berechneten geozentrischen
Örter der Reihe nach $M+ax+\beta y,M'+a'x+\beta'y,M''+a''x+\beta''y$ usw.
Nach den vorangegangenen Untersuchungen sind dann
$x,y$ so zu bestimmen, dass jene Größen mit den entsprechenden
$N,N',N''$ usw. möglichst gut übereinstimmen. Dabei ist die relative
Genauigkeit der Beobachtungen zu berücksichtigen. Das korrigierte
System von Elementen selbst lässt sich entweder auf gleiche Weise aus
$L,L'$ und den Entfernungen $D+x\delta,D'+y\delta'$ oder nach den
bekannten Regeln durch einfache Interpolation aus den ersten drei
Systemen von Elementen ableiten.
